\part{Интегралы}

\section{Определённый интеграл}

\begin{definition}
Пусть функция $f(x)$ определена на отрезке $[a,b]$.
Рассмотрим разбиение
\[
T:\quad a=x_0<x_1<\dots<x_n=b.
\]
Отрезки $[x_i,x_{i+1}]$ называются \emph{элементарными}.
Обозначим
\[
\Delta x_i=x_{i+1}-x_i,\qquad \lambda(T)=\max_{0\le i\le n-1}\Delta x_i.
\]
Для набора точек $\xi_i\in[x_i,x_{i+1}]$ сумма
\[
\sigma(T,\xi)=\sum_{i=0}^{n-1} f(\xi_i)\,\Delta x_i
\]
называется \emph{интегральной суммой Римана}.
\end{definition}

\begin{definition}
Если существует конечный предел
\[
\lim_{\lambda(T)\to 0}\sigma(T,\xi)=I,
\]
то функция $f$ называется \emph{интегрируемой} на $[a,b]$, а число $I$ — её \emph{определённым интегралом}:
\[
I=\int_a^b f(x)\,dx.
\]
Множество всех интегрируемых на $[a,b]$ функций обозначают $R[a,b]$.
\end{definition}

\begin{theorem}[Необходимое условие интегрируемости]
Если $f\in R[a,b]$, то $f$ ограничена на $[a,b]$.


\begin{proof}
Предположим противное: пусть функция $f$ не ограничена на $[a,b]$.
Тогда существует такой элементарный отрезок $[x_p,x_{p+1}]$, на котором $f$ неограничена.
Но тогда при выборе точки $\xi_p\in[x_p,x_{p+1}]$ значения $f(\xi_p)$ могут становиться сколь угодно большими, и соответствующие интегральные суммы не могут иметь конечного предела при $\lambda(T)\to 0$.
Противоречие.
\end{proof}

\end{theorem}

\section{Суммы Дарбу}

\begin{definition}
Пусть $f$ ограничена на $[a,b]$ и
\[
T:\quad a=x_0<x_1<\dots<x_n=b.
\]
Обозначим
\[
 m_i=\inf\limits_{x\in[x_i,x_{i+1}]} f(x),\qquad
 M_i=\sup\limits_{x\in[x_i,x_{i+1}]} f(x).
\]
Тогда \emph{нижняя} и \emph{верхняя суммы Дарбу} равны
\[
 s(T)=\sum_{i=0}^{n-1} m_i\,\Delta x_i,\qquad
 S(T)=\sum_{i=0}^{n-1} M_i\,\Delta x_i.
\]
\end{definition}

\begin{prop}
Для любой интегральной суммы $\sigma(T,\xi)$ выполнено
\[
 s(T)\le \sigma(T,\xi)\le S(T) \Rightarrow s(T) - \sup \land S(T) - \inf.
\]
\end{prop}

\begin{prop}
Если $T' = T\cup\{x\}$ —  $T$ без точки $x$, то
\[
S(T')\le S(T),\qquad s(T')\ge s(T).
\]
\end{prop}

\begin{prop}
Для любых разбиений $T$ и $T'$ выполняется
\[
 s(T)\le S(T').
\]
\end{prop}

Из предыдущего свойства следует, что множества всех нижних и верхних сумм ограничены:
\[
\boxed{s \le \underline I \le \overline I \le S} (*)
\]
\[
\underline I=\sup\{s(T)\},\qquad \overline I=\inf\{S(T)\}.
\]
Эти числа называются соответственно \emph{нижним} и \emph{верхним} интегралами.

\begin{theorem}[Теорема Дарбу]
Если $f$ ограничена на $[a,b]$ выполняется
\[
\lim_{\lambda(T)\to 0}s(T)=\underline I,\qquad
\lim_{\lambda(T)\to 0}S(T)=\overline I.
\]
\end{theorem}

\begin{remark}
Для краткости далее используется обозначение
\[
\omega_i=\sup\limits_{x\in[x_i,x_{i+1}]}f(x)-\inf\limits_{x\in[x_i,x_{i+1}]}f(x)
\]
— \emph{колебание} функции на элементарном отрезке.
Тогда
\[
S(T)-s(T)=\sum_{i=0}^{n-1}\omega_i\,\Delta x_i.
\]
\end{remark}

\section{Критерий интегрируемости функции}

\begin{theorem}
Для ограниченной на $[a,b]$ функции $f$ эквивалентны условия:
\[
 f\in R[a,b]\qquad\Longleftrightarrow\qquad \lim_{\lambda(T)\to 0}\bigl(S(T)-s(T)\bigr)=0.
\]


\begin{proof}
\textbf{Необходимость.}
Если $f\in R[a,b]$, то существует
\[
\lim_{\lambda(T)\to 0}\sigma(T,\xi)=I.
\]
По определению предела для любого $\varepsilon>0$ при достаточно малом $\lambda(T)$ имеем
\[
|\sigma(T,\xi)-I|<\frac{\varepsilon}{2}.
\]
Так как $s(T)\le \sigma(T,\xi)\le S(T)$, то
\[
I-\frac{\varepsilon}{2}\le s(T)\le S(T)\le I+\frac{\varepsilon}{2}.
\]
Следовательно,
\[
0\le S(T)-s(T)\le \varepsilon,
\]
значит $\lim_{\lambda(T)\to 0}(S(T)-s(T))=0$.

\medskip
\textbf{Достаточность.}
Пусть теперь
\[
\lim_{\lambda(T)\to 0}(S(T)-s(T))=0.
\]
Тогда по теореме Дарбу
\[
0\le \overline I-\underline I\le S(T)-s(T)\to 0,
\]
следовательно, $\underline I=\overline I=:I$.
Из свойства Дарбу получаем
\[
s(T)\le I\le S(T),
\]
а значит для любой интегральной суммы
\[
0\le |\sigma(T,\xi)-I|\le S(T)-s(T)\to 0.
\]
Следовательно, $\lim_{\lambda(T)\to 0}\sigma(T,\xi)=I$, то есть $f\in R[a,b]$.
\end{proof}

\end{theorem}

\begin{example}
Рассмотрим функцию Дирихле
\[
 f(x)=\begin{cases}
 1, & x\in\mathbb{Q},\\
 0, & x\in\mathbb{R}\setminus\mathbb{Q}.
 \end{cases}
\]
На любом отрезке $[a,b]$ и на любом элементарном отрезке рациональные и иррациональные числа всюду плотны, поэтому
\[
 m_i=0,\qquad M_i=1.
\]
Тогда
\[
 s(T)=0,\qquad S(T)=\sum_{i=0}^{n-1}\Delta x_i=b-a,
\]
то есть $S(T)-s(T)=b-a\neq 0$.
Значит, эта функция не интегрируема на $[a,b]$.
\end{example}

\section{Классы интегрируемых функций}

\begin{theorem}
Если $f\in C[a,b]$, то $f\in R[a,b]$.


\begin{proof}
По теореме Кантора функция, непрерывная на отрезке $[a,b]$, равномерно непрерывна.
Значит, для любого $\varepsilon>0$ существует такое $\delta>0$, что при $\lambda(T)<\delta$ для всех $i$ выполнено
\[
\omega_i<\frac{\varepsilon}{b-a}.
\]
Тогда
\[
0\le S(T)-s(T)=\sum_{i=0}^{n-1}\omega_i\,\Delta x_i
<\sum_{i=0}^{n-1}\frac{\varepsilon}{b-a}\,\Delta x_i
=\frac{\varepsilon}{b-a}\,(b-a)=\varepsilon.
\]
По критерию интегрируемости $f\in R[a,b]$.
\end{proof}
\end{theorem}

\begin{theorem}
Если $f$ монотонна на $[a,b]$, то $f\in R[a,b]$.

\begin{proof}
Пусть сначала $f$ неубывает на $[a,b]$.
Тогда на каждом элементарном отрезке $[x_i,x_{i+1}]$ имеем
\[
 m_i=f(x_i),\qquad M_i=f(x_{i+1}),
\]
поэтому
\[
\omega_i=M_i-m_i\le f(b)-f(a).
\]
Следовательно,
\[
S(T)-s(T)=\sum_{i=0}^{n-1}\omega_i\,\Delta x_i
\le \bigl(f(b)-f(a)\bigr)\sum_{i=0}^{n-1}\Delta x_i
=\bigl(f(b)-f(a)\bigr)(b-a).
\]
Для произвольного $\varepsilon>0$ выбираем разбиение с достаточно малым шагом $\lambda(T)$ так, чтобы сумма колебаний была меньше $\varepsilon$; тогда по критерию интегрируемости $f\in R[a,b]$.
Случай невозрастающей функции рассматривается аналогично.
\end{proof}
\end{theorem}
